\documentclass[11pt]{article}

\usepackage[margin=1in]{geometry}
\usepackage{amsmath,amssymb,bm}
\usepackage{booktabs}
\usepackage{array}
\usepackage{hyperref}
\usepackage{microtype}

\hypersetup{
  colorlinks=true,
  linkcolor=blue!45!black,
  urlcolor=blue!45!black,
  citecolor=blue!45!black
}

\newcommand{\PhiN}{\Phi}
\newcommand{\ind}{\mathbf{1}}
\newcolumntype{L}[1]{>{\raggedright\arraybackslash}p{#1}}

\title{\vspace{-2em}Signal-Support Simulation Design\\
\large Product-Mixture MAP versus Viroli-Style Probit Gibbs}
\author{}
\date{\vspace{-1.5em}}

\begin{document}
\maketitle

\section{Purpose}

The simulation asks when the proposed product-mixture MAP estimator can recover
a high-dimensional binary factor mixture model.  The working hypothesis is that
performance is governed less by $p$ alone and more by the strength and support
of the probit signal: how large the nonzero item loadings are, how many items
load on each factor, and how much signed cross-loading connects factors.

\section{Data-Generating Model}

Each replication generates
\[
  X_{ij} = \ind\{Z_{ij}>0\}, \qquad
  Z_{ij} = \alpha_j + \bm\lambda_j^\top \bm F_i + \epsilon_{ij},
  \qquad \epsilon_{ij}\sim N(0,1).
\]
The factor coordinates are independent marginal mixtures:
\[
  F_{ih}\sim \sum_{g=1}^{G_h}\pi_{hg}
  N(\mu_{hg},\sigma_{hg}^2), \qquad h=1,\ldots,H.
\]
The generated factors are standardized to the canonical marginal mean-zero,
unit-variance scale.  Item intercepts are included in all simulations using
the IFEval-like intercept generator.

\section{Simulation Grid}

\begin{center}
\begin{tabular}{L{0.34\linewidth}L{0.52\linewidth}}
\toprule
Quantity & Values \\
\midrule
Sample size & $n\in\{100,200\}$ \\
Items & $p\in\{500,1000,1500,2000\}$ for Product MAP \\
Gibbs item counts & $p\in\{500,1000\}$ only \\
Factors & $H\in\{5,10,15,20\}$ \\
Mixture components & $G_h=2$ for all $h$ or $G_h=3$ for all $h$ \\
Replications & 25 per setting \\
Mixture separation & 1 \\
Block sizes & balanced or IFEval-like unbalanced \\
Loading design & Cross/IFEval-like sparse signed cross-loadings \\
Loading strength & weak: $|\lambda|\sim U(1.25,1.75)$; strong: $|\lambda|\sim U(2.50,3.00)$ \\
Cross-loading probability & 0.075 or 0.20 \\
Cross-loading signs & random signs \\
\bottomrule
\end{tabular}
\end{center}

The loading strength range is applied to both primary loadings and nonzero
cross-loadings.  Primary loadings are positive.  Cross-loadings are randomly
positive or negative.  Representative loading matrices are generated by
\texttt{scripts/sample\_size/plot\_dgp\_loading\_heatmaps.R}.

\section{Product-Mixture MAP}

The proposed estimator has three stages.

\paragraph{Low-rank probit signal.}
First estimate the item intercepts and centered rank-$H$ probit signal
$B=F\Lambda^\top$ by maximizing
\[
  \ell(\bm\alpha,B)
  =
  \sum_{i,j}
  \left[
  X_{ij}\log \PhiN(\alpha_j+B_{ij})
  +(1-X_{ij})\log\{1-\PhiN(\alpha_j+B_{ij})\}
  \right],
  \qquad \mathrm{rank}(B)\le H.
\]
The implementation uses an EM-SVD update: compute conditional expectations of
the latent probit variables, update intercepts, and project the centered
working matrix onto rank $H$.  This stage estimates the signal subspace
directly and avoids repeated sampled-$Z$ augmentation in the proposed method.

\paragraph{Rotation.}
The SVD of $\widehat B$ gives left singular-vector scores spanning the signal
subspace, but the axes are rotationally unidentified.  The rotation stage
searches for an orthogonal rotation whose coordinates are well fit by
independent univariate Gaussian mixtures while also encouraging sparse
loadings.  The default implementation uses Riemannian rotation with an
$\ell_1$ loading penalty.

\paragraph{MAP refinement.}
Starting from the rotated scores, the refinement stage alternates updates of
factor scores, item intercepts/loadings, and marginal mixture parameters.  The
loadings are estimated with an $\ell_1$ penalty.  Refinement uses monotone
objective monitoring and stops only after the posterior objective is stable
and, by default, all marginal mixture updates have converged.

The default Product MAP penalties are
\[
  \lambda_{\mathrm{pretrain}}=
  \lambda_{\mathrm{rotation}}=
  \lambda_{\mathrm{refine}}=10.
\]

\section{Viroli-Style Gibbs Baselines}

Both Gibbs baselines fit the independent marginal mixture factor model with
probit augmentation.  Each iteration samples truncated latent probit variables,
factor scores, marginal component labels, mixture weights, mixture means,
mixture variances, item intercepts, and item loadings.  Posterior draws are
normalized to the same factor scale used in the DGP.

The Laplace baseline uses the Bayesian lasso scale-mixture representation for
the loading prior, with penalty 10.  The diffuse Gaussian baseline replaces
that prior with a weak Gaussian loading prior.  Both are run for 2000
iterations, discarding the first 1000 as burn-in.  Effective sample sizes are
recorded for parameters when enabled.

\section{Parallelization}

Product MAP uses one simulation chunk at a time by default, with internal
parallel workers for independent marginal mixture fits, itemwise loading
regressions, and subject-wise factor-score updates.  Gibbs is parallelized
across independent task chunks, so separate replications/configurations run
concurrently.  This avoids the overhead that previously made within-chain
parallel Gibbs unattractive in these R implementations.

\section{Recorded Results}

Each row records the full DGP settings, method settings, convergence
diagnostics, elapsed seconds, and recovery metrics:
\[
  \mathrm{RMSE}(F),\quad
  \mathrm{RMSE}(\Lambda),\quad
  \mathrm{RMSE}(\alpha),\quad
  \mathrm{RMSE}(\mu),\quad
  \mathrm{RMSE}(\sigma^2),\quad
  \mathrm{RMSE}(\pi).
\]
The product estimator also records first-stage signal diagnostics, including
relative Frobenius error for $\widehat B$ and the sine of the principal-angle
operator-norm error for the estimated signal subspace.  These diagnostics are
intended to separate failures of binary signal estimation from failures of
rotation or MAP refinement.

\section{Replication}

From the repository root:
\[
\texttt{Rscript scripts/sample\_size/run\_final\_product\_viroli\_simulation.R}.
\]
Progress plots can be regenerated at any point with
\[
\texttt{Rscript scripts/sample\_size/plot\_signal\_support\_simulation\_progress.R}.
\]
The full raw output is written to
\texttt{results/full/signal\_support\_grid}; selected GitHub-facing plots and
tables are written under \texttt{results/selected\_plots} and
\texttt{results/selected\_tables}.

\end{document}
